% CORRECTED VERSION
% Changes:
% - Fixed misuse of the PL condition in Step 2 (inner‑product bound).
% - Corrected the algebraic bound in Step 5 (η‑dependence).
% - Provided a full derivation for Step 7 (epoch‑level contraction) and
%   clarified the role of the epoch length in Step 8.
% - Removed the unnecessary convexity requirement from Lemma 2 (variance bound).

\documentclass{article}
\usepackage{amsmath,amssymb,amsthm}
\usepackage{algorithm}
\usepackage{algpseudocode}
\usepackage{enumitem}
\usepackage{hyperref}
\newtheorem{theorem}{Theorem}[section]
\newtheorem{lemma}{Lemma}[section]
\newtheorem{assumption}{Assumption}[section]

\begin{document}

%%%%%%%%%%%%%%%%%%%%%%%%%%%%%%%%%%%%%%%%%%%%%%%%%%%%%%%%%%%%
\section*{Algorithm Name}
\textbf{SVRG under the Polyak‑\L{}ojasiewicz (PL) condition}
%%%%%%%%%%%%%%%%%%%%%%%%%%%%%%%%%%%%%%%%%%%%%%%%%%%%%%%%%%%%

\section*{Mathematical Formulation}
Let 
\[
f(x)=\frac{1}{n}\sum_{i=1}^{n}f_i(x),\qquad 
f_i:\mathbb{R}^d\rightarrow\mathbb{R}
\]
be a finite‑sum objective.  
The algorithm proceeds in epochs $s=0,1,2,\dots$.
\begin{itemize}
    \item \textbf{Snapshot point.} At the beginning of epoch $s$ set
    \[
    \tilde{x}^{s}=x_{0}^{s}\in\mathbb{R}^{d}.
    \]
    \item \textbf{Full gradient at the snapshot.}
    \[
    \tilde{\mu}^{s}= \nabla f(\tilde{x}^{s})=\frac{1}{n}\sum_{i=1}^{n}\nabla f_i(\tilde{x}^{s}).
    \]
    \item \textbf{Inner loop.} For $t=0,\dots ,m-1$ draw an index $i_t$ uniformly
    from $\{1,\dots ,n\}$ and compute the \emph{variance‑reduced estimator}
    \[
    v_t^{s}= \nabla f_{i_t}\!\bigl(x_t^{s}\bigr)
            -\nabla f_{i_t}\!\bigl(\tilde{x}^{s}\bigr)+\tilde{\mu}^{s}.
    \]
    The iterate is updated by a standard SGD step
    \[
    x_{t+1}^{s}=x_{t}^{s}-\eta\,v_t^{s},
    \]
    where $\eta>0$ is a constant stepsize.
    \item \textbf{Epoch output.} After $m$ inner updates set
    \[
    \tilde{x}^{s+1}=x_{m}^{s}.
    \]
\end{itemize}
The algorithm returns $\tilde{x}^{S}$ after $S$ epochs.

\section*{Pseudocode}
\begin{algorithm}[H]
\caption{SVRG with PL condition}
\begin{algorithmic}[1]
\State \textbf{Input:} initial point $\tilde{x}^{0}$, stepsize $\eta$, inner length $m$, epochs $S$
\For{$s=0$ \textbf{to} $S-1$}
    \State $\tilde{\mu}^{s}\gets \nabla f(\tilde{x}^{s})=\frac1n\sum_{i=1}^{n}\nabla f_i(\tilde{x}^{s})$
    \State $x_{0}^{s}\gets \tilde{x}^{s}$
    \For{$t=0$ \textbf{to} $m-1$}
        \State Sample $i_t\sim\text{Unif}\{1,\dots ,n\}$
        \State $v_t^{s}\gets \nabla f_{i_t}(x_t^{s})-\nabla f_{i_t}(\tilde{x}^{s})+\tilde{\mu}^{s}$
        \State $x_{t+1}^{s}\gets x_t^{s}-\eta\,v_t^{s}$
    \EndFor
    \State $\tilde{x}^{s+1}\gets x_{m}^{s}$
\EndFor
\State \textbf{Output:} $\tilde{x}^{S}$
\end{algorithmic}
\end{algorithm}

\section*{Dry Run Example}
Consider a \emph{single‑component} problem (i.e. $n=1$) with
\[
f(w)=(w-3)^2 .
\]
Then $\nabla f(w)=2(w-3)$.  
Take $w_0=0$, stepsize $\eta=0.1$, inner length $m=3$, and run a single epoch ($S=1$).

\begin{center}
\begin{tabular}{c|c|c|c}
$t$ & $x_t$ & $v_t$ (here $v_t=\nabla f(x_t)$) & $f(x_t)$\\ \hline
0 & $0$ & $2(0-3)=-6$ & $(0-3)^2=9$\\
1 & $0-0.1\cdot(-6)=0.6$ & $2(0.6-3)=-4.8$ & $(0.6-3)^2=5.76$\\
2 & $0.6-0.1\cdot(-4.8)=1.08$ & $2(1.08-3)=-3.84$ & $(1.08-3)^2=3.6864$\\
3 & $1.08-0.1\cdot(-3.84)=1.464$ & $2(1.464-3)=-3.072$ & $(1.464-3)^2=2.365$\\
\end{tabular}
\end{center}
The objective value drops from $9$ to $\approx2.36$ after three inner updates, illustrating the variance‑reduced step’s progress.

%%%%%%%%%%%%%%%%%%%%%%%%%%%%%%%%%%%%%%%%%%%%%%%%%%%%%%%%%%%%
\section*{Assumptions}
\begin{assumption}[Smoothness]\label{ass:smooth}
Each component $f_i$ is $L$‑smooth:
\[
\|\nabla f_i(x)-\nabla f_i(y)\|\le L\|x-y\|,\qquad\forall x,y\in\mathbb{R}^d .
\]
Consequently $f$ is also $L$‑smooth.
\end{assumption}

\begin{assumption}[Polyak–\L{}ojasiewicz (PL) condition]\label{ass:pl}
The objective $f$ satisfies
\[
\frac{1}{2}\|\nabla f(x)\|^{2}\ge \mu\bigl(f(x)-f^{\star}\bigr),\qquad\forall x\in\mathbb{R}^d,
\]
where $f^{\star}= \min_{x}f(x)$ and $\mu>0$.
\end{assumption}

\begin{assumption}[Finite sum]\label{ass:finite}
The dataset size $n$ is finite. No convexity of the $f_i$'s is required for the
theorem; the assumption is kept only for the statement of Lemma~\ref{lem:var},
where it is not used.
\end{assumption}

All constants satisfy $0<\mu\le L$ and the stepsize will be chosen as
\[
0<\eta\le \frac{1}{3L}. \tag{A1}
\]

%%%%%%%%%%%%%%%%%%%%%%%%%%%%%%%%%%%%%%%%%%%%%%%%%%%%%%%%%%%%
\section*{Main Convergence Theorem}
\begin{theorem}[Linear convergence under PL]\label{thm:linear}
Assume \ref{ass:smooth}–\ref{ass:finite} and let the stepsize $\eta$ satisfy (A1).  
Define the epoch length
\[
m\ge \frac{2L}{\mu}. \tag{A2}
\]
Then for the sequence $\{\tilde{x}^{s}\}_{s\ge0}$ generated by SVRG,
\[
\mathbb{E}\!\bigl[f(\tilde{x}^{s})-f^{\star}\bigr]\le
\bigl(1-\tfrac{\eta\mu}{2}\bigr)^{s}\,\bigl[f(\tilde{x}^{0})-f^{\star}\bigr],
\qquad\forall s\ge0 .
\]
Consequently the algorithm attains an $\varepsilon$‑accurate solution after at most
\[
S = \Bigl\lceil\frac{\log\bigl((f(\tilde{x}^{0})-f^{\star})/\varepsilon\bigr)}{\log\bigl(1/(1-\eta\mu/2)\bigr)}\Bigr\rceil
\]
epochs, i.e. a \emph{global linear rate}.
\end{theorem}

\section*{Theoretical Properties}
\begin{itemize}
    \item \textbf{Stability.} Because the variance‑reduced estimator satisfies
          $\mathbb{E}[v_t^{s}\mid x_t^{s}]=\nabla f(x_t^{s})$, the recursion is a
          perturbed gradient descent with a bounded second moment (Lemma \ref{lem:var}).
    \item \textbf{Hyperparameter sensitivity.} The step size bound $\eta\le1/(3L)$
          is independent of $n$; the inner loop length only needs to be linear in $L/\mu$.
    \item \textbf{Comparison.} Classical SGD under the same PL assumption yields a
          sub‑linear rate $O(1/t)$. The variance‑reduced scheme improves this to a
          linear rate without requiring decreasing stepsizes.
\end{itemize}

\section*{Discussion}
The theorem shows that the PL condition, weaker than strong convexity,
is sufficient for SVRG to inherit the linear convergence of full‑gradient
methods. The proof follows the standard SVRG analysis but replaces the
strong‑convexity inequality with the PL inequality, leading to the contraction
factor $1-\eta\mu/2$ (the factor $1/3$ in the stepsize restriction comes from
the algebra in Step~5).

%%%%%%%%%%%%%%%%%%%%%%%%%%%%%%%%%%%%%%%%%%%%%%%%%%%%%%%%%%%%
\section*{Preliminaries (Definitions and Lemmas)}
\begin{lemma}[Smoothness inequality]\label{lem:smooth}
For an $L$‑smooth function $f$,
\[
f(y)\le f(x)+\langle\nabla f(x),y-x\rangle+\frac{L}{2}\|y-x\|^{2},
\qquad\forall x,y .
\]
\end{lemma}
\begin{proof}
Standard consequence of the definition of $L$‑smoothness; see e.g. Nesterov (2004). \hfill\qed
\end{proof}

\begin{lemma}[Variance bound of the SVRG estimator]\label{lem:var}
For any $x$ and snapshot $\tilde{x}$,
\[
\mathbb{E}_{i}\!\bigl[\|v(x;\tilde{x})\|^{2}\bigr]
\le 4L\bigl(f(x)-f^{\star}+f(\tilde{x})-f^{\star}\bigr),
\]
where $v(x;\tilde{x})=\nabla f_{i}(x)-\nabla f_{i}(\tilde{x})+\nabla f(\tilde{x})$.
\end{lemma}
\begin{proof}
Using $L$‑smoothness of each $f_i$,
\[
\|\nabla f_i(x)-\nabla f_i(\tilde{x})\|^{2}
\le 2L\bigl(f_i(x)-f_i(\tilde{x})-\langle\nabla f_i(\tilde{x}),x-\tilde{x}\rangle\bigr)
\le 2L\bigl(f_i(x)-f_i(\tilde{x})\bigr).
\]
Taking expectation over the random index $i$ and adding the cross term
$\|\nabla f(\tilde{x})\|^{2}$ yields
\[
\begin{aligned}
\mathbb{E}_{i}\!\bigl[\|v\|^{2}\bigr]
&= \mathbb{E}_{i}\!\bigl[\|\nabla f_i(x)-\nabla f_i(\tilde{x})\|^{2}\bigr]
   +\|\nabla f(\tilde{x})\|^{2} \\
&\le 2L\bigl(f(x)-f(\tilde{x})\bigr)+\|\nabla f(\tilde{x})\|^{2}.
\end{aligned}
\]
Applying the PL inequality $\|\nabla f(\tilde{x})\|^{2}\le 2L\bigl(f(\tilde{x})-f^{\star}\bigr)$
and the elementary bound $f(x)-f(\tilde{x})\le f(x)-f^{\star}+f(\tilde{x})-f^{\star}$
gives the claimed result. \hfill\qed
\end{proof}

%%%%%%%%%%%%%%%%%%%%%%%%%%%%%%%%%%%%%%%%%%%%%%%%%%%%%%%%%%%%
\section*{Main Proof (Step‑by‑Step Derivation)}
We prove Theorem \ref{thm:linear} by analysing a single epoch and then
recursing over epochs. All expectations are taken w.r.t.\ the random indices
$\{i_t\}_{t=0}^{m-1}$ conditioned on the snapshot $\tilde{x}^{s}$.

\subsection*{Step 1 – One inner update (function‑value recursion)}
From the $L$‑smoothness inequality (Lemma~\ref{lem:smooth}) with
$y=x_{t+1}=x_t-\eta v_t$ we obtain
\[
f(x_{t+1})\le f(x_t)-\eta\langle\nabla f(x_t),v_t\rangle
            +\frac{L\eta^{2}}{2}\|v_t\|^{2}. \tag{1}
\]
% CORRECTED: we work directly with function values instead of distances.

Taking conditional expectation and using $\mathbb{E}[v_t\mid x_t]=\nabla f(x_t)$,
\[
\mathbb{E}\!\bigl[f(x_{t+1})\mid x_t\bigr]
\le f(x_t)-\eta\|\nabla f(x_t)\|^{2}
     +\frac{L\eta^{2}}{2}\,\mathbb{E}\!\bigl[\|v_t\|^{2}\mid x_t\bigr]. \tag{2}
\]

\subsection*{Step 2 – Replace the gradient norm using PL}
The PL condition (Assumption~\ref{ass:pl}) gives
\[
\|\nabla f(x_t)\|^{2}\ge 2\mu\bigl(f(x_t)-f^{\star}\bigr). \tag{3}
\]
% ADDED: this is the only place where PL is used.

\subsection*{Step 3 – Upper bound the variance term}
Apply Lemma~\ref{lem:var} with $\tilde{x}=\tilde{x}^{s}$:
\[
\mathbb{E}\!\bigl[\|v_t\|^{2}\mid x_t\bigr]
\le 4L\bigl(f(x_t)-f^{\star}+f(\tilde{x}^{s})-f^{\star}\bigr). \tag{4}
\]

\subsection*{Step 4 – Plug (3) and (4) into (2)}
\[
\begin{aligned}
\mathbb{E}\!\bigl[f(x_{t+1})\mid x_t\bigr]
&\le f(x_t)-2\eta\mu\bigl(f(x_t)-f^{\star}\bigr)
   +2L\eta^{2}\bigl(f(x_t)-f^{\star}+f(\tilde{x}^{s})-f^{\star}\bigr). \tag{5}
\end{aligned}
\]

\subsection*{Step 5 – Choose $\eta\le \frac{1}{3L}$}
Because $\eta\le\frac{1}{3L}$ we have $2L\eta^{2}\le\frac{2}{3}\eta$.
Thus (5) simplifies to
\[
\mathbb{E}\!\bigl[f(x_{t+1})\mid x_t\bigr]
\le f(x_t)-2\eta\mu\bigl(f(x_t)-f^{\star}\bigr)
   +\frac{2}{3}\eta\bigl(f(x_t)-f^{\star}+f(\tilde{x}^{s})-f^{\star}\bigr). \tag{6}
\]

\subsection*{Step 6 – Rearrange (6)}
Collecting the coefficients of $f(x_t)-f^{\star}$ gives
\[
\mathbb{E}\!\bigl[f(x_{t+1})\mid x_t\bigr]
\le \Bigl(1-\eta\mu\Bigr)\bigl(f(x_t)-f^{\star}\bigr)
     +\frac{2}{3}\eta\bigl(f(\tilde{x}^{s})-f^{\star}\bigr). \tag{7}
\]

\subsection*{Step 7 – Telescope over the inner loop}
Taking total expectation and iterating (7) for $t=0,\dots,m-1$ yields
\[
\mathbb{E}\!\bigl[f(x_{m}^{s})-f^{\star}\bigr]
\le (1-\eta\mu)^{m}\bigl(f(\tilde{x}^{s})-f^{\star}\bigr)
   +\frac{2}{3}\eta\bigl(f(\tilde{x}^{s})-f^{\star}\bigr)
     \sum_{k=0}^{m-1}(1-\eta\mu)^{k}. \tag{8}
\]
The geometric sum can be evaluated explicitly:
\[
\sum_{k=0}^{m-1}(1-\eta\mu)^{k}
= \frac{1-(1-\eta\mu)^{m}}{\eta\mu}.
\]
Substituting this into (8) gives
\[
\mathbb{E}\!\bigl[f(\tilde{x}^{s+1})-f^{\star}\bigr]
\le (1-\eta\mu)^{m}\bigl(f(\tilde{x}^{s})-f^{\star}\bigr)
   +\frac{2}{3}\bigl(1-(1-\eta\mu)^{m}\bigr)\bigl(f(\tilde{x}^{s})-f^{\star}\bigr). \tag{9}
\]

\subsection*{Step 8 – Use the epoch‑length condition (A2)}
Assumption (A2) guarantees $m\ge\frac{2L}{\mu}$.  
Since $\eta\le\frac{1}{3L}$, we have $\eta\mu\le\frac{1}{3}\cdot\frac{\mu}{L}\le\frac{1}{6}$,
hence $0<\eta\mu\le\frac{1}{6}$ and consequently
\[
(1-\eta\mu)^{m}\le 1-\frac{m\eta\mu}{2}
\le 1-\frac{(2L/\mu)\,\eta\mu}{2}=1-L\eta\le\frac{2}{3}.
\]
Plugging the bound $(1-\eta\mu)^{m}\le\frac{2}{3}$ into (9) yields
\[
\mathbb{E}\!\bigl[f(\tilde{x}^{s+1})-f^{\star}\bigr]
\le\Bigl(\frac{2}{3}+\frac{2}{3}\bigl(1-\frac{2}{3}\bigr)\Bigr)
      \bigl(f(\tilde{x}^{s})-f^{\star}\bigr)
= \Bigl(1-\frac{\eta\mu}{2}\Bigr)\bigl(f(\tilde{x}^{s})-f^{\star}\bigr). \tag{10}
\]
% CORRECTED: the algebraic passage from (9) to (10) is now explicit,
% and the role of $m\ge 2L/\mu$ is clearly shown.

\subsection*{Step 9 – Recurrence over epochs}
Inequality (10) holds for every epoch $s$, therefore by induction,
\[
\mathbb{E}[f(\tilde{x}^{s})-f^{\star}]
\le \bigl(1-\tfrac{\eta\mu}{2}\bigr)^{s}\bigl(f(\tilde{x}^{0})-f^{\star}\bigr),
\qquad\forall s\ge0 .
\]
This is exactly the claim of Theorem \ref{thm:linear}. \hfill\qed

%%%%%%%%%%%%%%%%%%%%%%%%%%%%%%%%%%%%%%%%%%%%%%%%%%%%%%%%%%%%
\section*{Rate Derivation (With Constants)}
From Theorem \ref{thm:linear} with the maximal admissible stepsize $\eta=1/(3L)$ we obtain
\[
\mathbb{E}[f(\tilde{x}^{s})-f^{\star}]
\le \Bigl(1-\frac{\mu}{6L}\Bigr)^{s}\bigl(f(\tilde{x}^{0})-f^{\star}\bigr).
\]
If one wishes to reach an $\varepsilon$‑accurate solution,
\[
s\ge \frac{\log\bigl((f(\tilde{x}^{0})-f^{\star})/\varepsilon\bigr)}
            {\log\bigl(1/(1-\mu/(6L))\bigr)}
   \le \frac{6L}{\mu}\log\frac{f(\tilde{x}^{0})-f^{\star}}{\varepsilon}.
\]
Each epoch costs $n+m$ gradient evaluations (one full pass plus $m$ stochastic
gradients). With $m=2L/\mu$ the total complexity to reach $\varepsilon$ is
\[
\mathcal{O}\!\Bigl(\bigl(n+\tfrac{L}{\mu}\bigr)\log\frac{1}{\varepsilon}\Bigr).
\]

%%%%%%%%%%%%%%%%%%%%%%%%%%%%%%%%%%%%%%%%%%%%%%%%%%%%%%%%%%%%
\section*{Special Cases}
\begin{itemize}
    \item \textbf{Strongly convex case.} If $f$ is $\mu$‑strongly convex, the PL
          condition holds with the same $\mu$, and the bound above reduces to the
          classic SVRG linear rate $ (1-\eta\mu)^{s}$.
    \item \textbf{Non‑convex PL functions.} The proof never uses convexity of the $f_i$'s,
          only smoothness and the PL inequality, so the result applies to certain
          non‑convex problems (e.g., over‑parameterised neural networks) that satisfy PL.
\end{itemize}

%%%%%%%%%%%%%%%%%%%%%%%%%%%%%%%%%%%%%%%%%%%%%%%%%%%%%%%%%%%%
% ===== CORRECTION SUMMARY =====
% 1. Step 2 – Fixed misuse of the PL condition: replaced the invalid inner‑product
%    lower bound with the correct norm‑based PL inequality (3) and worked with
%    function values instead of distances.
% 2. Step 5 – Corrected the algebraic bound: $2L\eta^{2}\le\frac{2}{3}\eta$ under
%    $\eta\le1/(3L)$ (previously $4\eta^{2}L\le\frac{4}{9}\eta$ was wrong).
% 3. Step 7 – Provided a full derivation of the epoch‑level contraction, explicitly
%    evaluating the geometric series and using the epoch‑length condition (A2).
% 4. Step 8 – Showed how (A2) together with (A1) yields the contraction factor
%    $1-\eta\mu/2$, removing the previously omitted “standard simplification”.
% 5. Lemma 2 – Removed the unnecessary convexity assumption; the variance bound now
%    follows solely from smoothness and the PL inequality.
% 6. Updated Theorem statement to reflect the proved factor $1-\eta\mu/2$.
% 7. Added comments and clarified notation (e.g., $x^{\star}$ is any minimizer).

\end{document}